Considera la funció
\[
f(x)=x^3-2x+1 .
\]

\begin{apartats}

\apartat[1,25]{1}
Troba l'equació de la recta tangent a la gràfica de $f$ en el punt d'abscissa $x=1$.

\begin{solucio}
$f(1)=1-2+1=0$ i $f'(x)=3x^2-2$, d'on $f'(1)=1$.\\
Tangent: $y-f(1)=f'(1)(x-1)$, és a dir $y=x-1$.
\end{solucio}

\begin{nomesllarg}
\apartat{0,75}
Troba l'equació de la recta normal a la gràfica de $f$ en el mateix punt.

\begin{solucio}
La normal és perpendicular a la tangent: el seu pendent és $-\dfrac{1}{f'(1)}=-1$.\\
Normal: $y-0=-(x-1)$, és a dir $y=-x+1$.
\end{solucio}
\end{nomesllarg}

\begin{tria}{tangent-tasca}
\itemtria{original}{0,75}{1,25}
La gràfica de $g(x)=x^2+ax-4$ passa pel punt $P(2,6)$. Troba $a$ i l'equació de la recta
tangent a la gràfica de $g$ en $P$.

\begin{solucio}
$g(2)=4+2a-4=2a=6$, d'on $\boxed{a=3}$.\\
$g'(x)=2x+3$ i $g'(2)=7$. Tangent: $y-6=7(x-2)$, és a dir $y=7x-8$.
\end{solucio}

\itemtria{tangents-simetriques}{0,75}{1,25}
Troba les rectes tangents a la gràfica de $f$ en els punts d'abscissa $x=2$ i $x=-2$, i comprova
que són paral·leles. Passa el mateix en $x=a$ i $x=-a$, sigui quin sigui $a$? Justifica-ho.

\begin{solucio}
$f'(x)=3x^2-2$, i per tant $f'(2)=f'(-2)=10$. Com que $f(2)=5$ i $f(-2)=-3$, les tangents són
$y=10x-15$ i $y=10x+17$: tenen el mateix pendent i no són la mateixa recta, i per tant són
paral·leles.\\
Sí, passa sempre: $f'(-a)=3(-a)^2-2=3a^2-2=f'(a)$. Les tangents en $x=a$ i $x=-a$ tenen el
mateix pendent, i són paral·leles (o la mateixa recta, si $a=0$).
\end{solucio}
\end{tria}

\end{apartats}
